\documentclass[11pt]{article}

\usepackage{amsmath,amssymb,mathtools,bm}
\usepackage[margin=1in]{geometry}

\newcommand{\E}{\mathbb{E}}
\newcommand{\R}{\mathbb{R}}
\newcommand{\N}{\mathcal{N}}
\newcommand{\ELBO}{\mathrm{ELBO}}
\newcommand{\tr}{\mathrm{tr}}
\setlength{\parindent}{0pt}

\title{20260320 Number Crunch}
\author{}
\date{}

\begin{document}
\maketitle

\section*{Background and Problem Intuition}

\textbf{Mathematical process:} Start from Bayes' rule, then convert inference into optimisation.
\begin{align}
p(\bm\theta\mid\bm y)
&= \frac{p(\bm y\mid\bm\theta)p(\bm\theta)}{p(\bm y)}, \\
p(\bm y)
&= \int p(\bm y,\bm\theta)\,d\bm\theta,
\end{align}
where $p(\bm y)$ is typically intractable.

Introduce a tractable variational density $q(\bm\theta)$ and write
\begin{align}
\log p(\bm y)
&= \log p(\bm y)\int q(\bm\theta)\,d\bm\theta \\
&= \int q(\bm\theta)\log p(\bm y)\,d\bm\theta \\
&= \int q(\bm\theta)f(\bm\theta)\,d\bm\theta,
\quad f(\bm\theta)\equiv \log p(\bm y) \\
&= \E_q\!\left[f(\bm\theta)\right] \\
&= \E_q\!\left[\log p(\bm y)\right] \\
&= \E_q\!\left[\log \frac{p(\bm y,\bm\theta)}{p(\bm\theta\mid\bm y)}\right] \\
&= \E_q\!\left[\log \left(\frac{p(\bm y,\bm\theta)}{q(\bm\theta)}\cdot
\frac{q(\bm\theta)}{p(\bm\theta\mid\bm y)}\right)\right] \\
&= \E_q\!\left[\log \frac{p(\bm y,\bm\theta)}{q(\bm\theta)}\right]
 +\E_q\!\left[\log \frac{q(\bm\theta)}{p(\bm\theta\mid\bm y)}\right] \\
&= \E_q[\log p(\bm y,\bm\theta)]-\E_q[\log q(\bm\theta)]
 +\E_q\!\left[\log \frac{q(\bm\theta)}{p(\bm\theta\mid\bm y)}\right].
\end{align}

\textbf{Sequence of events (plain-language):}
\begin{enumerate}
\item Start from $\log p(\bm y)$, which is constant with respect to $\bm\theta$.
\item Multiply by $1$ using normalisation of $q$: $\int q(\bm\theta)\,d\bm\theta=1$.
\item Rewrite as an integral: $\log p(\bm y)=\int q(\bm\theta)\log p(\bm y)\,d\bm\theta$.
\item Recognise expectation form: $\int q(\bm\theta)f(\bm\theta)\,d\bm\theta=\E_q[f(\bm\theta)]$.
\item Use Bayes identity inside the logarithm: $p(\bm y)=\frac{p(\bm y,\bm\theta)}{p(\bm\theta\mid\bm y)}$.
\item Insert $q(\bm\theta)/q(\bm\theta)$ to create ELBO and KL terms.
\item Split product and quotient with log rules, then split expectations linearly.
\item Identify
\[\ELBO(q)=\E_q[\log p(\bm y,\bm\theta)]-\E_q[\log q(\bm\theta)]\]
and
\[\mathrm{KL}\!\left(q(\bm\theta)\,\|\,p(\bm\theta\mid\bm y)\right)=\E_q\!\left[\log\frac{q(\bm\theta)}{p(\bm\theta\mid\bm y)}\right].\]
\end{enumerate}

Define
\begin{align}
\ELBO(q)
&= \E_q[\log p(\bm y,\bm\theta)]-\E_q[\log q(\bm\theta)] \\
&= \E_q[\log p(\bm y,\bm\theta)] + H(q),
\end{align}
so that
\begin{equation}
\log p(\bm y) = \ELBO(q) + \mathrm{KL}\!\left(q(\bm\theta)\,\|\,p(\bm\theta\mid\bm y)\right).
\end{equation}
Since $\mathrm{KL}(\cdot\|\cdot)\ge 0$, maximising $\ELBO(q)$ is equivalent to minimising divergence to the true posterior.

\section*{Key Notation}


\section*{Optimisation Problem Formulation}

\subsection*{Decision Variables}

\textbf{Mathematical process:} Parameter counting and feasible-domain specification.
\begin{align}
\bm\phi &= (\bm\mu_\beta,\bm\Sigma_\beta,a_e,b_e),\\
\bm\mu_\beta &\in \R^p,\quad \bm\Sigma_\beta\in\mathbb S_{++}^p,\quad a_e>0,\ b_e>0,\\
d &= p + \frac{p(p+1)}{2} + 2.
\end{align}
For $p=3$: $d=3+6+2=11$.

\subsection*{Objective Function}

\textbf{Mathematical process:} Expand $\ELBO(\bm\phi)$ term-by-term under mean-field
$q(\bm\beta,\tau_e)=q(\bm\beta)\,q(\tau_e)$ with
$q(\bm\beta)=\N(\bm\mu_\beta,\bm\Sigma_\beta)$ and
$q(\tau_e)=\mathrm{Gamma}(a_e,b_e)$ (shape/rate parameterisation).

Write the five-term decomposition:
\begin{equation}
\ELBO(\bm\phi)= T_1+T_2+T_3+T_4+T_5,
\end{equation}
where
\[
T_1=\E_q[\log p(\bm y\mid\bm\beta,\tau_e)],\quad
T_2=\E_q[\log p(\bm\beta)],\quad
T_3=\E_q[\log p(\tau_e)],\quad
T_4=-\E_q[\log q(\bm\beta)],\quad
T_5=-\E_q[\log q(\tau_e)].
\]

\textbf{Core building-block expectations} (used throughout):
\begin{align}
\E_q[\tau_e] &= \frac{a_e}{b_e}, &
\E_q[\log\tau_e] &= \psi(a_e)-\log b_e,\\
\E_q\!\left[\|\bm y-\bm X\bm\beta\|^2\right]
  &= \|\bm y-\bm X\bm\mu_\beta\|^2+\tr(\bm X^\top\bm X\bm\Sigma_\beta), &
\E_q[\bm\beta^\top\bm\beta]
  &= \bm\mu_\beta^\top\bm\mu_\beta+\tr(\bm\Sigma_\beta).
\end{align}
The first two follow from Gamma moment formulae; the third uses
$\E_q[\bm\beta\bm\beta^\top]=\bm\Sigma_\beta+\bm\mu_\beta\bm\mu_\beta^\top$
and the trace–quadratic identity; the fourth is the same identity with $\bm A=\bm I$.

\subsubsection*{$T_1$: Expected log-likelihood}

From $\bm y\mid\bm\beta,\tau_e\sim\N(\bm X\bm\beta,\tau_e^{-1}\bm I_n)$:
\begin{equation}
\log p(\bm y\mid\bm\beta,\tau_e)
= \frac{n}{2}\log\tau_e - \frac{n}{2}\log(2\pi)
  - \frac{\tau_e}{2}\|\bm y-\bm X\bm\beta\|^2.
\end{equation}
Taking the expectation and using mean-field independence ($\tau_e\perp\bm\beta$ under $q$):
\begin{equation}
T_1 = \frac{n}{2}\bigl(\psi(a_e)-\log b_e\bigr)
    - \frac{n}{2}\log(2\pi)
    - \frac{a_e}{2b_e}\Bigl(\|\bm y-\bm X\bm\mu_\beta\|^2
      +\tr(\bm X^\top\bm X\bm\Sigma_\beta)\Bigr).
\end{equation}

\subsubsection*{$T_2$: Expected log-prior on $\bm\beta$}

From $\bm\beta\sim\N(\bm 0,\lambda_\beta^{-1}\bm I_p)$:
\begin{equation}
\log p(\bm\beta)
= \frac{p}{2}\log\lambda_\beta - \frac{p}{2}\log(2\pi)
  - \frac{\lambda_\beta}{2}\bm\beta^\top\bm\beta,
\end{equation}
\begin{equation}
T_2 = \frac{p}{2}\log\lambda_\beta - \frac{p}{2}\log(2\pi)
    - \frac{\lambda_\beta}{2}\Bigl(\bm\mu_\beta^\top\bm\mu_\beta
      +\tr(\bm\Sigma_\beta)\Bigr).
\end{equation}

\subsubsection*{$T_3$: Expected log-prior on $\tau_e$}

From $\tau_e\sim\mathrm{Gamma}(\alpha_e,\gamma_e)$ (shape/rate):
\begin{equation}
\log p(\tau_e)
= \alpha_e\log\gamma_e - \log\Gamma(\alpha_e)
  + (\alpha_e-1)\log\tau_e - \gamma_e\tau_e,
\end{equation}
\begin{equation}
T_3 = \alpha_e\log\gamma_e - \log\Gamma(\alpha_e)
    + (\alpha_e-1)\bigl(\psi(a_e)-\log b_e\bigr)
    - \gamma_e\,\frac{a_e}{b_e}.
\end{equation}

\subsubsection*{$T_4$: Entropy of $q(\bm\beta)$}

The differential entropy of a $p$-variate Gaussian is
\begin{equation}
T_4 = H\!\bigl(q(\bm\beta)\bigr)
= \frac{p}{2}\bigl(1+\log(2\pi)\bigr)+\frac{1}{2}\log\lvert\bm\Sigma_\beta\rvert.
\end{equation}
With the Cholesky factor $\bm\Sigma_\beta=\bm L\bm L^\top$ this becomes
$\log\lvert\bm\Sigma_\beta\rvert=2\sum_i\log L_{ii}$, which is numerically stable.

\subsubsection*{$T_5$: Entropy of $q(\tau_e)$}

The differential entropy of a $\mathrm{Gamma}(a,b)$ (shape/rate) distribution is
\begin{equation}
T_5 = H\!\bigl(q(\tau_e)\bigr)
= a_e - \log b_e + \log\Gamma(a_e) + (1-a_e)\psi(a_e).
\end{equation}

\subsubsection*{Full ELBO (up to constants in $\bm\phi$)}

Dropping all terms that depend only on hyperparameters $(\lambda_\beta,\alpha_e,\gamma_e)$:
\begin{align}
\ELBO(\bm\phi)
&= \frac{n}{2}(\psi(a_e)-\log b_e)
   -\frac{a_e}{2b_e}\Bigl(\|\bm y-\bm X\bm\mu_\beta\|^2
   +\tr(\bm X^\top\bm X\bm\Sigma_\beta)\Bigr)
   \tag{$T_1$}\\
&\quad
   -\frac{\lambda_\beta}{2}\bigl(\bm\mu_\beta^\top\bm\mu_\beta+\tr(\bm\Sigma_\beta)\bigr)
   \tag{$T_2$}\\
&\quad
   +(\alpha_e-1)(\psi(a_e)-\log b_e)
   -\gamma_e\frac{a_e}{b_e}
   \tag{$T_3$}\\
&\quad
   +\tfrac{1}{2}\log\lvert\bm\Sigma_\beta\rvert
   \tag{$T_4$}\\
&\quad
   +\log\Gamma(a_e)+(1-a_e)\psi(a_e)-\log b_e+a_e
   +C,
   \tag{$T_5$}
\end{align}
where $C$ absorbs all terms independent of $\bm\phi=(\bm\mu_\beta,\bm\Sigma_\beta,a_e,b_e)$.

\subsection*{Constraints}

\textbf{Mathematical process:} Convert constrained variables to unconstrained parameters.
\begin{align}
\bm\Sigma_\beta &= \bm L\bm L^\top,\quad \bm L \text{ lower triangular with } L_{ii}>0,\\
a_e &= e^{\eta_a},\quad b_e=e^{\eta_b},\quad \eta_a,\eta_b\in\R.
\end{align}
This guarantees $\bm\Sigma_\beta\succ0$ and $a_e,b_e>0$ by construction.

\subsection*{Computational Characteristics}

\textbf{Mathematical process:} Leading-order complexity accounting.
\begin{align}
\text{Gradient evaluation} &:\ O(np^2),\\
\text{Newton step (Hessian build + solve)} &:\ O(np^2+d^3),\\
\text{BFGS update} &:\ O(d^2).
\end{align}

\section*{Test Case: Bayesian Linear Regression}

\textbf{Mathematical process:} Specify generative model, variational family, and joint log-density components.
\begin{align}
\bm y\mid\bm\beta,\tau_e &\sim \N(\bm X\bm\beta,\tau_e^{-1}\bm I_n),\\
\bm\beta &\sim \N(\bm 0,\lambda_\beta^{-1}\bm I_p),\\
\tau_e &\sim \mathrm{Gamma}(\alpha_e,\gamma_e).
\end{align}

Useful joint log form (up to constants):
\begin{align}
\log p(\bm y,\bm\beta,\tau_e)
&= \frac{n}{2}\log\tau_e
-\frac{\tau_e}{2}\|\bm y-\bm X\bm\beta\|^2
-\frac{\lambda_\beta}{2}\bm\beta^\top\bm\beta \\
&\quad +(\alpha_e-1)\log\tau_e-\gamma_e\tau_e + C.
\end{align}

\section*{Overview of Optimisation Methods}

\subsection*{Methods considered}


\subsection*{Baseline Method: CAVI\footnote{Blei, Kucukelbir, and McAuliffe (2017), Variational Inference: A Review for Statisticians (Journal of the American Statistical Association).}}

\textbf{Mathematical process:} Coordinate-wise maximisation of ELBO.
\begin{align}
\bm\Sigma_\beta &\leftarrow \left(\frac{a_e}{b_e}\bm X^\top\bm X\right)^{-1},\\
\bm\mu_\beta &\leftarrow \frac{a_e}{b_e}\,\bm\Sigma_\beta\bm X^\top\bm y,\\
a_e &\leftarrow \alpha_e+\frac n2,\\
b_e &\leftarrow \gamma_e+\frac12\Big(\|\bm y-\bm X\bm\mu_\beta\|^2+\tr(\bm X^\top\bm X\bm\Sigma_\beta)\Big).
\end{align}
Each block update solves
\begin{equation}
q_j^{(t+1)} = \arg\max_{q_j}\ \ELBO(q_j,q_{-j}^{(t)}),
\end{equation}
so ELBO is non-decreasing across iterations.

\subsection*{Mapping ENEL445 Methods to VI Problem}

\textbf{Mathematical process:} Express optimisation updates on $\bm\phi$.

\textbf{Gradient ascent}
\begin{equation}
\bm\phi^{(k+1)} = \bm\phi^{(k)} + \alpha_k\nabla\ELBO(\bm\phi^{(k)}).
\end{equation}

\textbf{Newton's method}
\begin{equation}
\bm\phi^{(k+1)}
= \bm\phi^{(k)} + \alpha_k\big[\nabla^2\ELBO(\bm\phi^{(k)})\big]^{-1}\nabla\ELBO(\bm\phi^{(k)}).
\end{equation}

\textbf{BFGS (inverse form)}
\begin{align}
\bm s_k &= \bm\phi_{k+1}-\bm\phi_k,
\qquad
\bm y_k = \nabla\ELBO(\bm\phi_{k+1})-\nabla\ELBO(\bm\phi_k),\\
\rho_k &= (\bm y_k^\top\bm s_k)^{-1},\\
\bm B_{k+1} &= (\bm I-\rho_k\bm s_k\bm y_k^\top)\bm B_k(\bm I-\rho_k\bm y_k\bm s_k^\top)+\rho_k\bm s_k\bm s_k^\top.
\end{align}

\section*{Solution Plan}

\subsection*{Phase 1: Data Generation and Baseline}

\textbf{Mathematical process:} Derive and verify all analytical quantities required before coding.
\begin{align}
\text{Deliverables}:
\quad &\ELBO(\bm\phi),\ \nabla\ELBO(\bm\phi),\ \nabla^2\ELBO(\bm\phi),\\
&\text{CAVI updates, and numerical checks against finite differences.}
\end{align}

\subsection*{Phase 2: Implementation (in progress)}


\subsection*{Success Metrics}

\textbf{Mathematical process:} Define quantitative stopping and agreement criteria.
\begin{align}
\text{Function tolerance:}
\quad &|\ELBO^{(k+1)}-\ELBO^{(k)}|<10^{-6},\\
\text{Gradient tolerance:}
\quad &\|\nabla\ELBO(\bm\phi^{(k)})\|_2<10^{-6},\\
\text{Method agreement:}
\quad &|\ELBO_{\text{method}}^\star-\ELBO_{\text{CAVI}}^\star|\le \varepsilon.
\end{align}

\section*{Progress To Date}

\subsection*{Work in progress using an iterative / Agile method}

\textbf{Mathematical process identified as drafted/partial:}
\begin{align}
&\text{(i) ELBO derivation from model assumptions,}\\
&\text{(ii) CAVI closed-form updates,}\\
&\text{(iii) analytical gradients for } d=11,\\
&\text{(iv) Hessian block structure,}\\
&\text{(v) preliminary BFGS and penalty formulations.}
\end{align}

\subsection*{Preliminary Results}

\textbf{Mathematical process:} Local structure checks before implementation.
\begin{align}
\nabla_{\bm\mu_\beta}\ELBO
&\propto \bm X^\top(\bm y-\bm X\bm\mu_\beta),\\
\text{Hessian block sign checks}
&\Rightarrow \text{local concavity in selected directions},\\
 d&=11\ \Rightarrow\ \text{Newton linear-algebra cost is tractable at this scale.}
\end{align}

\subsection*{Next Steps}


\section*{Alignment with Course Requirements}


\section*{Timeline}


\section*{Summary}


\end{document}
